\section{Weighted-defect form of the global curvature question}
\label{sec:weightedcurvature}
There is a particularly direct equivalent form of the
remaining monotonicity condition, using the same variable
$t$ as the earlier Csordas--Dimitrov differential operator.

\begin{proposition}[Exact weighted-defect equivalence]
Let $s:(0,\infty)\to(0,\infty)$ be three times
continuously differentiable and put
\[
\Phi(u)=s(u^2),\quad V(u)=-\log\Phi(u),\quad
q(t)=\frac{\mathcal Ts(t)}{s(t)^2}
=-(\log s)''(t).
\]
For every $u>0$,
\begin{align}
V'''(u)&=12u\,q(u^2)+8u^3q'(u^2)\nonumber\\
&=8\left.\frac{d}{dt}
[t^{3/2}q(t)]\right|_{t=u^2}.
\label{eq:weightedeq}
\end{align}
Consequently
\[
V'''(u)\ge0\ \ (u>0)
\quad\Longleftrightarrow\quad
t^{3/2}\frac{\mathcal Ts(t)}{s(t)^2}
\text{ is nondecreasing on }t>0.
\]
\end{proposition}
\begin{proof}
Set $L(t)=\log s(t)$. Chain differentiation
gives
\[
V''(u)=-2L'(u^2)-4u^2L''(u^2)
\]
and hence
\[
V'''(u)=-12uL''(u^2)-8u^3L'''(u^2).
\]
Substituting $q=-L''$ proves the first identity,
and differentiating $t^{3/2}q(t)$ proves
the second. Since the change $u\mapsto u^2$
preserves the positive half-line, the final
equivalence follows.
\end{proof}

The known second-level theorem controls the
\emph{sign} of $\mathcal Ts$ but does not by
itself control the derivative of the weighted
quotient. Thus the new proof target is
\[
3q(t)+2tq'(t)\ge0,\qquad t>0.
\]
It is a one-dimensional inequality, involving
no Fourier integration, and could be tackled
through Jacobi's differential equation,
a compact interval certificate, or a direct
comparison inequality for $\mathcal Ts$ and
$s$. Even full success establishes only a
common finite Laguerre positivity band, not RH.
